\song{Pasážová revolta (Karel Kryl)}
%Mara_2018_06_15

\ve{}
\ch{Emi}Nosíme z módy \ch{Ami}kopretiny, \ch{Emi}{č}ímž \ch{H}okrádáme \ch{Emi H7}stáda\\
a \ch{Emi}vůl, kdys' jméno \ch{Ami}obětiny, \ch{Emi}je titul \ch{H7}kama\ch{Emi}ráda,\\
\ch{Ami}na obou nohách vietnamku \ch{Emi}a jako komfort \ch{H7}hlavu,\\
\ch{Emi}na klopě placku \ch{Ami}jak psí známku - \ch{Emi}znak \ch{H7}příslušnosti k \ch{Emi}davu.\\
I naše \ch{Ami}generace má svoje \ch{D}prominenty,\\
program je \ch{G}rezignace \ch{Ami}a facky \ch{H7 (H7, H)}argumenty,\\
potlesk je k \ch{Emi}umlčení a pískot \ch{C}na pochvalu\\
a místo \ch{D}přesvědčení jen pití \ch{H}piva - z žalu.

\ve{}
Pod zadkem stránku Dikobrazu, vzýváme zlatý tele,\\
sedáme v koutcích u obrazů, čekáme Spasitele,\\
civíme lačně na měďáky, my - Gottwaldovi vnuci,\\
a nadáváme na měšťáky, tvoříce Revoluci.\\
I naše generace má svoje kajícníky\\
a fízly z honorace a skromné úředníky\\
a tvory bez svědomí a plazy bez páteře\\
a život v bezvědomí a lásku - k nedůvěře.

\ve{}
Už nejsme, nejsme to, co kdysi, už známe ohnout záda,\\
umíme dělat komromisy a zradit kamaráda\\
a, vděčni dnešní realitě, líbáme cizí ruce\\
a jednou zajdem na úbytě z tý smutný revoluce.\\
I v naší generaci už máme pamětníky\\
a vlastní emigraci a vlastní mučedníky.

\re{Rec.:}
\textit{\ch{Emi}A s hubou rozmlácenou dnes \ch{H7}zůstali jsme \ch{Emi}němí,\\
ne, nejsme na kolenou - ryjeme \ch{H7}držkou v \ch{Emi}zemi!}

\esong